% Part: incompleteness
% Chapter: incompleteness-provability
% Section: tarski-thm

\documentclass[../../../include/open-logic-section]{subfiles}

\begin{document}

\olfileid{inc}{inp}{tar}

\olsection{સત્યની અવ્યાખ્યેયતા}

\emph{વ્યાખ્યેયતા}નો ખ્યાલ અંકગણિતની ભાષા માટે ઔપચારિક અર્થવિજ્ઞાન હોવા પર આધાર રાખે છે. આપણે અંકગણિતની ભાષામાં સૂત્રો અને વાક્યોનો ગણ વર્ણવ્યો છે. ``અભિપ્રેત અર્થઘટન'' મુજબ એ વાક્યો કુદરતી સંખ્યાઓ વિશે દાવા કરે છે; આવા દાવા સાચા કે ખોટા હોઈ શકે. $\Struct{N}$ એ એવો !!{structure} હોય જેનો પ્રદેશ $\Nat$ છે અને જેમાં અંકગણિતની ભાષાના ચિહ્નોનું પ્રમાણભૂત અર્થઘટન છે. ત્યારે $\Sat{N}{!A}$નો અર્થ ``પ્રમાણભૂત અર્થઘટનમાં $!A$ સાચું છે'' થાય છે.

\begin{defn}
કુદરતી સંખ્યાઓનો સંબંધ $R(x_1,\dots,x_k)$, $\Struct{N}$માં \emph{વ્યાખ્યેય} ત્યારે અને માત્ર ત્યારે જ છે જ્યારે અંકગણિતની ભાષામાં એવું સૂત્ર $!A(x_1,\dots,x_k)$ હોય કે બધાં $n_1,\dots,n_k$ માટે $R(n_1,\dots,n_k)$ ત્યારે અને માત્ર ત્યારે જ હોય જ્યારે $\Sat{N}{!A(\num n_1,\dots,\num n_k)}$ હોય.
\end{defn}

બીજી રીતે કહીએ તો, સંબંધ $\Struct{N}$માં વ્યાખ્યેય ત્યારે અને માત્ર ત્યારે જ છે જ્યારે તે સિદ્ધાંત $\Th{TA}$માં નિરૂપણક્ષમ હોય. અહીં $\Th{TA} = \Setabs{!A}{\Sat{N}{!A}}$ અંકગણિતનાં સાચાં વાક્યોનો ગણ છે. (આ વાત તરત સ્પષ્ટ ન લાગે તો વ્યાખ્યાઓ ફરી જોઈને પોતાને ખાતરી કરાવો.)

\begin{lem}
દરેક સંગણનીય સંબંધ $\Struct{N}$માં વ્યાખ્યેય છે.
\end{lem}

\begin{proof}
$\Th{Q}$માં કોઈ સંબંધને નિરૂપતું સૂત્ર $\Struct{N}$માં એ જ સંબંધ વ્યાખ્યાયિત કરે છે એ ચકાસવું સહેલું છે.
\end{proof}

હવે વિપરીત દિશા પણ સાચી છે કે નહીં એવું પૂછીએ: શું $\Struct{N}$માં વ્યાખ્યેય દરેક સંબંધ સંગણનીય છે? જવાબ ના છે. દાખલા તરીકે:

\begin{lem}
અટકવાનો સંબંધ $\Struct{N}$માં વ્યાખ્યેય છે.
\end{lem}

\begin{proof}
યાદ કરો કે ક્લીનીનું સામાન્ય સ્વરૂપ પ્રમેય કહે છે: દરેક આંશિક સંગણનીય વિધેય~$f$ માટે એવો સૂચક~$e$ છે કે બધાં $x \in \Nat$ માટે $f(x) = U(\umin{s}{T(e,x,s)})$; અહીં $U$ અને $T$ આદિમ પુનરાવર્તી અને તેથી સર્વત્ર વ્યાખ્યાયિત છે. આથી $f(x)$ વ્યાખ્યાયિત હોય (એટલે સંગણના અટકે) ત્યારે અને માત્ર ત્યારે જ કોઈ~$s$ માટે $T(e,x,s)$ સાચું હોય.

હવે $H$ અટકવાનો સંબંધ હોય, એટલે
\[
H = \Setabs{\tuple{e,x}}{\lexists[s][T(e, x, s)]}.
\]
$!D_T$ એ $\Struct{N}$માં $T$ને વ્યાખ્યાયિત કરતું સૂત્ર હોય. તો
\[
H = \Setabs{\tuple{e,x}}{\Sat{N}{\lexists[s][!D_T(\num e, \num x, s)]}},
\]
અને તેથી $\lexists[s][!D_T(z, x, s)]$ એ $\Struct{N}$માં $H$ને વ્યાખ્યાયિત કરે છે.
\end{proof}

\begin{prob}
બતાવો કે $Q(n) \defiff n \in \Setabs{\Gn{!A}}{\Th{Q} \Proves !A}$ અંકગણિતમાં વ્યાખ્યેય છે.
\end{prob}

પોતે $\Th{TA}$ વિશે શું? શું તે અંકગણિતમાં વ્યાખ્યેય છે? એટલે, શું ગણ $\Setabs{\Gn{!A}}{\Sat{N}{!A}}$ અંકગણિતમાં વ્યાખ્યેય છે? ટાર્સ્કીનું પ્રમેય તેનો નકારાત્મક જવાબ આપે છે.

\begin{thm}
\ollabel{thm:tarski}
અંકગણિતનાં સાચાં !!{sentence}નો ગણ અંકગણિતમાં વ્યાખ્યેય નથી.
\end{thm}

\begin{proof}
માનો કે $!D(x)$ તેને વ્યાખ્યાયિત કરે છે, એટલે $\Sat{N}{!A}$ ત્યારે અને માત્ર ત્યારે જ જ્યારે $\Sat{N}{!D(\gn{!A})}$. નિશ્ચિત-બિંદુ લેમ્માથી એવું સૂત્ર $!A$ છે કે $\Th{Q} \Proves !A \liff \lnot !D(\gn{!A})$ અને તેથી $\Sat{N}{!A \liff \lnot !D(\gn{!A})}$. પરંતુ ત્યાર પછી $\Sat{N}{!A}$ ત્યારે અને માત્ર ત્યારે જ જ્યારે $\Sat{N}{\lnot !D(\gn{!A})}$; આ વાત $!D(y)$ અંકગણિતનાં સાચાં વિધાનોનો ગણ વ્યાખ્યાયિત કરે છે એવી ધારણાને વિરોધે છે.
\end{proof}

ટાર્સ્કીએ આ વિશ્લેષણને સત્યના વધુ વ્યાપક દાર્શનિક ખ્યાલ પર લાગુ કર્યું. કોઈપણ ભાષા $L$ માટે તેમણે દલીલ કરી કે $L$ માટે સત્યનો યોગ્ય ખ્યાલ દરેક વાક્ય $X$ માટે આ શરત સંતોષે:
\begin{quote}
`$X$' ત્યારે અને માત્ર ત્યારે જ સાચું છે જ્યારે $X$.
\end{quote}
અંગ્રેજી માટેનું ટાર્સ્કીનું વારંવાર ઉલ્લેખાતું ઉદાહરણ ગુજરાતી અર્થમાં આ છે:
\begin{quote}
`બરફ સફેદ છે' ત્યારે અને માત્ર ત્યારે જ સાચું છે જ્યારે બરફ સફેદ છે.
\end{quote}
પરંતુ કર્ણીય વિધેયનું નિરૂપણ કરી શકે એટલી શક્તિશાળી કોઈપણ ભાષા અને કોઈપણ ભાષાકીય વિધેયધર્મ $T(x)$ માટે આપણે એવું વાક્ય $X$ રચી શકીએ જે ``$X$ ત્યારે અને માત્ર ત્યારે જ જ્યારે $T(\text{`$X$'})$ નથી'' સંતોષે. સત્યનો વિધેયધર્મ કોઈ વાક્યને એકસાથે સાચું અને ખોટું ન ઠરાવે એવી આપણી માગણી હોવાથી ટાર્સ્કીએ તારવ્યું કે ભાષાની મર્યાદાની બહાર કોઈ રીતે ગયા વિના તે ભાષાનાં બધાં વાક્યો માટે સત્યનો વિધેયધર્મ નિર્દેશી શકાતો નથી. બીજા શબ્દોમાં, ભાષા માટેનો સત્યનો વિધેયધર્મ એ ભાષામાં જ વ્યાખ્યાયિત થઈ શકતો નથી.

\end{document}
